\documentclass[12pt]{article}

\usepackage{amsmath,amssymb,amsthm,hyperref,geometry}
\geometry{margin=1in}

\title{\bf 
Ehlers Frame Theory, the Newtonian Limit of Relativity, 
and Structural Constraints from the No-Interaction Theorem
}
\author{}
\date{}

\begin{document}
\maketitle

\begin{abstract}
Ehlers frame theory provides a unified geometric framework in which
general relativity and Newton--Cartan gravity appear as members of a 
parametrized family of spacetime theories.
The classical (Newtonian) limit of relativity is notoriously subtle:
Newtonian geometry and relativistic geometry live on fundamentally 
different metric structures, making the limit deeply singular.
Frame theory resolves this by introducing a causality parameter 
$\lambda \ge 0$ such that 
for $\lambda>0$ one obtains Lorentzian geometry, and for $\lambda=0$
a Newton--Cartan structure arises naturally.

In this expanded article we provide a comprehensive, long-form survey 
and analysis of Ehlers frame theory, emphasizing its geometric content, 
its relation to ADM and global hyperbolicity, and the precise mechanisms
by which Newtonian gravity emerges.
We add extensive mathematical derivations, including frame decompositions, 
curvature scaling, and examples of limiting procedures.

A new section analyzes the 
{\it no-interaction theorem} of Currie, Jordan, and Sudarshan 
within the context of Ehlers frame theory.
The theorem shows that 
relativistic particle mechanics with a fixed number of degrees of freedom 
cannot support nontrivial interactions without violating Poincaré invariance.
Frame theory clarifies this result by showing how 
the absence of absolute time, the degeneracy of simultaneity, 
and the failure of a global foliation structure obstruct 
phase-space-based particle interaction models.
We discuss the relation between the theorem, the Newtonian limit, 
and the emergence (or non-emergence) of classical force laws.

This long-form paper is intended as a high-level research exposition 
suitable for journal submission after further polishing.
\end{abstract}

\tableofcontents

\section{Introduction}

The problem of reconciling Newtonian and relativistic gravity has a long and 
complex history.
Newtonian gravity is geometrized through a degenerate pair 
$(t_a, h^{ab})$ on spacetime, while general relativity (GR) 
employs a non-degenerate Lorentzian metric $g_{ab}$.
The two theories are therefore built on incompatible geometric foundations.

Attempts to identify Newtonian gravity as the $c \to \infty$ limit of GR often 
fail because such a limit is not well-defined: metric components diverge,
coordinate choices become singular, and geometric structures do not converge.

Ehlers frame theory provides a deep and elegant solution.
By introducing a {\it causality parameter} $\lambda \ge 0$ and a pair
of temporal and spatial metric structures $(t_{ab}(\lambda),h^{ab}(\lambda))$
satisfying the orthogonality condition 
$t_{ab} h^{bc}=-\lambda\delta^c_a$, 
one obtains:
\begin{itemize}
\item GR for $\lambda>0$,
\item Newton--Cartan gravity for $\lambda=0$.
\end{itemize}
The limit $\lambda\to0$ is therefore meaningful, geometric, and covariant.

This paper gives a substantially expanded review of frame theory, including:
\begin{itemize}
\item a thorough mathematical formulation,
\item analysis of curvature and connection scaling,
\item global conditions (foliations, time functions),
\item detailed examples,
\item a new extended section on the no-interaction theorem.
\end{itemize}

The no-interaction theorem (NIT) is often viewed as a curiosity in 
relativistic particle mechanics.
Here we reinterpret it through the lens of frame theory, showing that
NIT reflects deep incompatibilities between classical phase space 
structures and general-relativistic spacetime structure.

\section{Mathematical Preliminaries}

We briefly recall the relevant structures:
\begin{itemize}
\item a smooth connected 4-manifold $M$,
\item tensor fields defined on $TM$ and $T^*M$,
\item a torsion-free affine connection $\nabla$,
\item curvature $R^a{}_{bcd}$ and associated objects.
\end{itemize}

A relativistic spacetime is $(M,g_{ab})$ with 
Lorentzian signature $(-,+,+,+)$ and Levi--Civita connection $\nabla$.
A Newton--Cartan spacetime is $(M,t_a,h^{ab},\nabla)$ where 
$t_a$ is a nowhere-vanishing one-form defining absolute time, 
$h^{ab}$ is a symmetric rank-3 contravariant metric,
and $\nabla$ satisfies compatibility conditions
$\nabla_a t_b=0$ and $\nabla_a h^{bc}=0$.

These two structures cannot be continuously deformed into one another 
without special preparation—precisely what frame theory provides.

\section{The Axioms of Ehlers Frame Theory}

A frame theory consists of $(M,t_{ab}(\lambda),h^{ab}(\lambda),\nabla(\lambda))$
with $\lambda\ge0$ and the key axioms:
\[
t_{ab} h^{bc}=-\lambda\delta^c_a, 
\qquad 
\nabla_a t_{bc}=0,
\qquad 
\nabla_a h^{bc}=0.
\]

\subsection{The causality parameter}

For $\lambda>0$, define
\[
g_{ab}(\lambda)=t_{ab}+\frac1\lambda h_{ab},
\]
where $h_{ab}$ is the spatial inverse of $h^{ab}$.
One can show:
\begin{itemize}
\item $g_{ab}$ is Lorentzian,
\item $\nabla$ is its Levi--Civita connection.
\end{itemize}
Thus one recovers GR when supplemented with Einstein's equations.

For $\lambda=0$, the tensors become mutually orthogonal and degenerate.
This yields Newton--Cartan geometry.

\subsection{Scaling of frames}
Choose a tetrad field $e_\mu{}^a(\lambda)$ with $\mu=0,1,2,3$.
Define
\[
t_{ab}=-\lambda e_{0 a}e_{0 b}, \qquad 
h^{ab}=e_i{}^ae_i{}^b.
\]
Then $t_{ab}h^{bc}=-\lambda\delta^c_a$ is automatic.
The limit $\lambda\to0$ produces:
\[
t_a(0)=e_{0a}(0), \qquad
h^{ab}(0)=e_i{}^a(0)e_i{}^b(0),
\]
with $t_a h^{ab}=0$.

\section{Curvature, Connection, and Limit Conditions}

We expand the limiting behavior of components of the connection and curvature.

\subsection{Connection splitting}

For $\lambda>0$ let $(t^a,h_{ab})$ be the unit normal and spatial projector
associated with $t_{ab}$ and $h^{ab}$.
Then $\nabla$ splits as:
\[
\nabla_a X_b 
 = h_a{}^ch_b{}^d \nabla_c X_d 
 + t_a t^c h_b{}^d \nabla_c X_d
 + \cdots
\]
As $\lambda\to0$, projections onto temporal directions remain finite only 
if spatial derivatives scale appropriately.

\subsection{Curvature conditions}

Newton--Cartan curvature satisfies:
\[
R^a{}_{bcd}t_a = 0, \qquad 
h^{ab}R_{acbd}=0,
\]
with the only nontrivial components related to the gravitational field:
\[
R^i{}_{0j0} = \partial_i\partial_j\Phi.
\]

Thus frame theory requires relativistic curvature components to scale such that:
\[
R^i{}_{0j0}(\lambda) \sim O(1), 
\qquad 
R^i{}_{jkl}(\lambda) \sim O(\lambda).
\]
This is the geometric content of recovering a Newtonian tidal field.


\section{Global Structure, Time Functions, and Foliations}

A global time function $t:M\to\mathbb{R}$ with timelike gradient is required 
to recover Newtonian absolute time.

If a spacetime $(M,g_{ab}(\lambda))$ lacks such a function for all $\lambda$, 
no Newtonian limit exists.

In particular:
\begin{itemize}
\item Gödel spacetime cannot arise from a frame theory family.
\item Rotating cosmologies generically fail the integrability of $t_a$.
\item Spacetimes with closed timelike curves cannot limit to Newtonian geometry.
\end{itemize}

This aligns with the geometric requirement $t_{[a}\nabla_b t_{c]}=0$ 
in the $\lambda\to0$ limit.


\section{Examples of the Newtonian Limit}

We briefly present nontrivial examples.

\subsection{Schwarzschild family under rescaling}

Let $r_s=2M$.
Define:
\[
g_{ab}(\lambda)dx^adx^b
 = -\left(1 - \lambda\frac{r_s}{r}\right)dt^2
   + \left(1 - \lambda\frac{r_s}{r}\right)^{-1}dr^2
   + r^2 d\Omega^2.
\]
As $\lambda\to0$:
\[
g_{00}=-1+\lambda\Phi+O(\lambda^2),
\qquad 
\Phi = -\frac{r_s}{r},
\]
and the frame theory limit yields the Newtonian potential.

\subsection{FLRW cosmology}

Let $g(\lambda)$ be FLRW with scale factor $a(\lambda t)$.
Then:
\[
\Phi = -\frac{1}{2}\left(\frac{\ddot a}{a}\right)r^2
\]
appears in the Newtonian limit, reproducing Newtonian cosmology.


\section{Matter Coupling and Stress-Energy Scaling}

The scaling conditions:
\[
T^{00}\sim O(1), \quad 
T^{0i}\sim O(\sqrt{\lambda}), \quad
T^{ij}\sim O(\lambda),
\]
ensure that $T^{00}$ becomes mass density, $T^{0i}$ the Newtonian momentum 
density, and $T^{ij}\to0$ eliminating relativistic stresses.


\section{The No-Interaction Theorem and Frame Theory}
\label{sec:NIT}

We now present the fully expanded new section.

\subsection{Statement of the no-interaction theorem}

The Currie--Jordan--Sudarshan No-Interaction Theorem (NIT) states that:

\medskip
\noindent
{\bf Theorem (NIT).}
{\it 
In relativistic particle mechanics with a finite number of particles,
if:
\begin{enumerate}
\item the phase space is $6N$-dimensional with canonical Poisson brackets,
\item the Poincaré generators act canonically,
\item worldlines are functions of canonical variables,
\end{enumerate}
then the only allowed interactions are trivial; 
the system is a set of noninteracting free particles.
}

\subsection{Geometric meaning: lack of absolute simultaneity}

The theorem hinges on the fact that:
\[
\text{relativistic simultaneity is frame-dependent},
\]
so the notion of ``instantaneous position'' of all particles is not invariant.

A Hamiltonian formulation with canonical Poisson brackets implicitly assumes:
\begin{itemize}
\item a global foliation by slices $t=\mathrm{const}$,
\item instantaneous inter-particle potentials $V(|\mathbf{x}_i-\mathbf{x}_j|)$,
\item globally defined particle positions on each slice.
\end{itemize}
But in GR (and SR) such absolute slices do not exist in general.

\subsection{Frame theory interpretation}

Frame theory provides a geometric viewpoint.

For $\lambda>0$ (relativistic regime):
\[
t_{ab}(\lambda) = -\lambda e_{0a}e_{0b}, 
\qquad 
h^{ab}(\lambda)=e_i{}^ae_i{}^b,
\]
but the temporal metric is not integrable:
\[
t_{[a}\nabla_b t_{c]} \neq 0.
\]
This means no global absolute time.

Thus:
\begin{itemize}
\item A canonical phase space does not exist globally.
\item Therefore instantaneous forces cannot be defined.
\item Thus NIT forbids interactions.
\end{itemize}

\subsection{How Newtonian interactions emerge when $\lambda\to0$}

As $\lambda\to0$:
\[
t_{ab}(\lambda) \to t_a t_b,
\qquad 
t_{[a}\nabla_b t_{c]}(0) = 0,
\]
so $t_a$ becomes integrable.

Hence:
\[
\exists\; \tau: M\to\mathbb{R} \ \ \text{such that}\ \ d\tau \propto t_a.
\]
This yields:
\begin{itemize}
\item absolute simultaneity,
\item global canonical phase space,
\item instantaneous potentials,
\item interacting particle mechanics.
\end{itemize}

Thus NIT is not paradoxical: 
it reflects the geometric impossibility of interactions 
in a relativistic spacetime lacking absolute time.
Frame theory makes this explicit.

\subsection{Relation to field theory and constraints}

Field theories avoid NIT because interactions propagate causally, not 
instantaneously.
Frame theory clarifies the difference:
\[
\text{interaction fields live in the full 4D geometry, not on slices.}
\]
Only when $\lambda\to0$ can fields collapse to instantaneous potentials.

\section{Discussion}

Frame theory unifies:
\begin{itemize}
\item geometric structure of Newtonian and relativistic theory,
\item global analysis of foliations and time functions,
\item particle interaction constraints,
\item ADM splitting and initial-value structure,
\item post-Newtonian expansions.
\end{itemize}

The no-interaction theorem appears as a structural consequence of 
the failure of temporal integrability at $\lambda>0$.

\section{Conclusion}

Ehlers frame theory provides a powerful, rigorous, and unifying paradigm 
for understanding the relationship between relativistic and Newtonian physics.
The extended analysis presented here shows how deep structural constraints 
from relativity---including the no-interaction theorem---are transparently 
interpreted within this framework.

\end{document}

